% LSPV-2 REPORT::TAGDILUTE1::2026-10-12T21:25:36Z
% !TEX program = xelatex
\documentclass[11pt]{article}
\usepackage{fontspec}
\usepackage[margin=0.86in]{geometry}
\usepackage{amsmath,amssymb,amsthm,booktabs,longtable,tabularx,enumitem,xurl,hyperref,microtype}
\hypersetup{colorlinks=true,linkcolor=blue,urlcolor=blue,citecolor=blue}
\newtheorem{theorem}{Theorem}
\newtheorem{definition}{Definition}
\newcommand{\Core}{\mathcal C_{2,2}^{\pm}}
\newcommand{\EXACT}{\textbf{PROVED/EXACT}}
\newcommand{\KNOWN}{\textbf{KNOWN THEOREM}}
\newcommand{\COMP}{\textbf{COMPUTATIONAL}}
\newcommand{\CAND}{\textbf{CANDIDATE}}
\newcommand{\OPEN}{\textbf{OPEN}}
\newcommand{\RET}{\textbf{FALSIFIED/RETIRED}}
\newcommand{\same}{\textsc{same\_realization}}
\title{TAGDILUTE1: Raw encoding density is not a semantic invariant}
\author{LSPV-2 / UNIVERSAL\_CORE\_TESTS}
\date{12 October 2026, 21:25 UTC}
\begin{document}
\maketitle

\begin{center}
\fbox{\parbox{0.94\linewidth}{
\textbf{Executive status.} P versus NP remains \OPEN.  Even when a violation
predicate depends only on the decoded clean-core instance and is therefore
constant on every representation fiber, uniform density over raw encodings
can be diluted by an arbitrary factor through unequal redundant fibers.
For \(M\) semantic instances, one bad instance, and tag multiplicity
\(q=2^r\), its density can be changed from \(1/M\) to
\(1/(1+(M-1)q)\) without changing any decoded instance or decision bit.
Thus \texttt{RAW-ENCODING-DENSITY-AS-UNIVERSAL-BRIDGE1} is \RET{}.
A fixed canonical encoding, semantic pushforward distribution, or proved
bounded-fiber sampler remains \OPEN, as does the direct whole-execution lower
bound.}}
\end{center}

\section{Claim card, predecessor, and selected obstacles}
\begin{description}[leftmargin=0cm,labelwidth=0cm]
\item[Claim ID.] \texttt{REDUNDANT-ENCODING-DENSITY-DILUTION1}.
\item[Status before run.] \OPEN{} repair proposed after DENSNOGO1: require a
mutation-invariant property and specify a distribution.
\item[Exact prefix.] For all integers \(M\ge2,r\ge0\), there exists a
polynomial-time recognizable redundant encoding of an \(M\)-object semantic
layer and a representation-invariant singleton bad set whose raw-uniform
density is \(1/(1+(M-1)2^r)\).
\item[Model.] Self-delimiting bit strings; a uniform deterministic parser
validates and removes the tag in linear time, then invokes an arbitrary total
decider on the canonical instance. Parsing, normalization, decision,
extraction, reconstruction, and the one-bit output are charged.
\item[\same{} objects.] One raw input string, its decoded clean-core formula,
the normalizing wrapper execution, and the base-decider execution it invokes.
\item[What would falsify it.] A counting error, failure of uniform
recognizability, or dependence of the bad predicate on the redundant tag.
\end{description}

The cyclic predecessor is LSPV-1. Its report
\path{research/agents/lspv-1/2026-10-13_0608_LSPV-1_DENSNOGO1.tex}
was read back from canonical GitHub with header
\texttt{LSPV-1 REPORT::DENSNOGO1} and the final verification sentinel. No
recovery request was needed. The two active obstacles are (i) defining a
density that is covariant under polynomial-time redundant encodings, and
(ii) transferring any resulting violation mass to unavoidable total cost for
every arbitrary uniform SAT decider in the same ordinary-machine realization.

\section{Semantic clean-core layer and two encoding schemes}
Adopt the explicit DENSNOGO1 slice.  Both six-variable blocks have positive
supports
\[
123,\quad145,\quad246,\quad356
\]
and negative supports
\[
N_0=\{135,146,236,245\},\qquad
N_1=\{136,145,235,246\}.
\]
Let \(G_0,G_1\) be the corresponding signed formulas and set
\[
 F_u=G_{u_1}\sqcup\cdots\sqcup G_{u_k},\qquad
 D_k=\{F_u:u\in\{0,1\}^k\}.                              \tag{1}
\]
Both blocks are duplicate-free proper signed Monotone
\(3\)SAT-\((2,2)\) instances, and both are satisfiable.  Thus
\(D_k\subseteq\Core\), \(M_k=|D_k|=2^k\), and its canonical indexed length is
\(N_k=\Theta(k\log k)\). Let \(z_k=F_{0^k}\).

For a fixed \(r\), put \(T_r=\{0,1\}^r\).  A self-delimiting raw input is a
pair \((F_u,t)\), and the decoder discards \(t\).  Consider:
\begin{align}
 E^{\mathrm{bal}}_{k,r}&=D_k\times T_r,                    \tag{2}\\
 E^{\mathrm{dil}}_{k,r}&=\{(z_k,0^r)\}
   \cup\{(F_u,t):u\ne0^k,\ t\in T_r\}.                  \tag{3}
\end{align}
Both are recognizable in \(O(N_k+r)\) bit operations.  All raw strings have
the same length \(N'_k=N_k+r+O(\log r)\).  The second grammar is deliberately
hostile but uniform: compare \(u\) with \(0^k\), and if equal require
\(t=0^r\).  No SAT call is needed for validation.

Let the bad predicate be semantic:
\[
 B(F_u,t)=1\quad\Longleftrightarrow\quad F_u=z_k.          \tag{4}
\]
It is constant on every decoder fiber and invariant under all permutations,
Gray recodings, or XOR masks of the valid tag set.  Thus this is stronger than
a one-presentation trace mutation.

\section{Exact density-dilution theorem}
\begin{theorem}[\texttt{REDUNDANT-ENCODING-DENSITY-DILUTION1}]
For the uniform measures on (2) and (3),
\[
 \mu_{\rm bal}(B)=\frac1{M_k},\qquad
 \mu_{\rm dil}(B)=\frac1{1+(M_k-1)2^r}.                 \tag{5}
\]
Consequently
\[
 \frac{\mu_{\rm bal}(B)}{\mu_{\rm dil}(B)}
 =\frac{1+(M_k-1)2^r}{M_k}.                              \tag{6}
\]
The decoded semantic layer and the bad semantic object are identical in the
two schemes.  For \(r=k\), the raw density changes from \(2^{-k}\) to
\(1/(1+(2^k-1)2^k)=2^{-2k+o(1)}\); for \(r=N_k\), the change is exponential
in the encoded length while normalization remains linear.
\end{theorem}

\begin{proof}
In the balanced scheme every semantic object has a fiber of size \(2^r\), so
the numerator is \(2^r\) and the support has size \(M_k2^r\).  In the diluted
scheme the bad fiber has size one and every other fiber has size \(2^r\), so
the support has size \(1+(M_k-1)2^r\).  Equation (4) proves representation
invariance.  The parser described above proves uniform recognizability and
the length/cost claims.  Division gives (5)--(6). \qed
\end{proof}

More generally, for a decoder \(\pi:E\to D\), semantic bad set \(S\subset D\),
and uniform raw measure,
\[
 \Pr_{e\in E}[\pi(e)\in S]
 =\frac{\sum_{x\in S}|\pi^{-1}(x)|}
        {\sum_{x\in D}|\pi^{-1}(x)|}.                    \tag{7}
\]
Thus raw-uniform density is a fiber-weighted semantic distribution, not a
semantic invariant. This is a complete counting theorem, independent of any
finite experiment.

\section{Algorithm, costs, and scope}
For any total deterministic decider \(A\) on canonical formulas, define
\(W_A(F_u,t)=A(F_u)\) after validation and tag removal. Then
\[
 T_{W_A}(N'_k)=O(N'_k)+T_A(N_k),\qquad
 S_{W_A}(N'_k)=O(N'_k)+S_A(N_k),                          \tag{8}
\]
under a direct buffer implementation; a streaming parser can reduce the tag
workspace. Compatibility extraction and witness reconstruction are zero, and
the output is one bit. DIRECT, COPY, DEAD, DELAY, INTERP, and RECOMP wrappers
can be added without changing (5) or the answer.

The theorem does \emph{not} refute a density statement under one permanently
fixed canonical encoding. It refutes only the use of uniform raw-encoding
density as an encoding-agnostic or semantics-only bridge. The smallest repair
must do at least one of the following:
\begin{enumerate}
\item fix a canonical encoding and charge canonicalization and sampling;
\item specify a semantic distribution \(\nu\) and sample raw encodings by
first sampling \(F\sim\nu\), preventing fiber multiplicity from changing
semantic mass; or
\item prove polynomial upper and lower bounds on all relevant fiber ratios
under every admitted encoding transformation.
\end{enumerate}
Even after that repair, one must prove inverse-polynomial violation mass for
\emph{every} arbitrary decider and transfer detection to a model-matched
whole-execution lower bound. Those arrows remain \OPEN.

\section{Bounded reproducible experiment}
The script
\path{experiments/review/redundant_encoding_density_dilution_audit.py}
writes
\path{results/raw/review/2026-10-13_redundant_encoding_density_dilution.json}.
One invocation completed in 0.0384 seconds; no chunking was needed. It checked
all \((k,r)\in[1,8]\times[0,8]\), both exact collision counts, four bijective
tag recodings, and 60,000 seeded sampling trials (seed 202610130625). It also
verified the two clean blocks have respectively 20 and 18 witnesses and ran
864 answer-preservation checks across clean-core, 2SAT, Horn, affine,
bounded-treewidth, independent-product, padding, redundant, and
alternate-encoding controls under all six wrappers. All 61,487 assertions
passed.

Script SHA-256:
\path{22e3c44d8cd1d7ae44fa215e2a71f1458f36d6f931be670825b491ba58df432c}.
JSON SHA-256:
\path{b7840eb77846b03c7183c93a5976a31ad4d2f47852ac9ad78edd7b19a182172e}.
The finite audit is \COMP{} and validates implementation only. The theorem is
the symbolic fiber count above.

\section{Literature and barrier-transfer preflight}
Goldreich--Goldwasser--Ron define property testing relative to a specified
object representation, access model, distance, and proximity parameter
\cite{GGR98}.  It is a \textbf{DIRECTLY TRANSFERABLE THEOREM} only after those
objects are fixed; it does not assert invariance under redundant encodings.
Goldreich--Ron explicitly define both an object distance and a distribution
distance for distributions of huge objects~\cite{GR23}; this is a
\textbf{USEFUL STRUCTURAL ANALOGY REQUIRING A NEW TRANSFER THEOREM}. Equation
(7) identifies the pushforward measure that a SAT density claim would need,
but their testing results do not provide an all-decider SAT lower bound.
Current primary work separates query complexity from computational complexity
in property testing~\cite{FPR26}; therefore a tester theorem alone cannot be
promoted to ordinary SAT running time.

Coding-theoretic puncturing, shortening, and quotient measures are
\textbf{USEFUL STRUCTURAL ANALOGIES}: coordinate multiplicities must be
transported with the measure. VLSI routing, communication cuts, branching
programs, and proof systems are \textbf{FAILED TRANSFERS} without a compiler
that maps the same arbitrary ordinary execution and preserves every cost.

The density-dilution theorem relativizes and algebrizes. It defines no large
constructive circuit property and supplies no Natural-Proofs escape. Hence it
is a route retirement, not a proof technique for P versus NP. Baker--Gill--
Solovay and Aaronson--Wigderson remain barriers, never premises
\cite{BGS75,AW09}.

\section{Post-quantum technique ledger}
For the diluted decoder, lossiness/collision entropy is an
\textbf{EXACT OBJECT/THEOREM TRANSFER}:
\[
 |E|=1+(M-1)q,qquad
 \operatorname{Coll}(\pi)=1+(M-1)q^2,qquad q=2^r.        \tag{9}
\]
The exact Cauchy--Schwarz bound
\(|\operatorname{im}\pi|\ge |E|^2/\operatorname{Coll}(\pi)\) detects a large
semantic image but does not bound decision time. The fiber imbalance is
algorithmically exploitable: the parser simply discards the tag.

Worst-case/average-case architecture is an \textbf{ANALOGY/CANDIDATE}: a
repair needs an unconditional sampler whose pushforward mass is controlled
for every arbitrary solver and whose total extraction cost is charged.
Special-soundness/collision extraction, Schur squares, quadratic relations,
Hilbert functions, syzygies, and graded Betti numbers are
\textbf{FAILED TRANSFER} for the present one-bit interface. Hash or
random-oracle normalization is inadmissible without removing the assumption.

\section{CORPUS REVIEWED}
This run surveyed the current GitHub tree and read root \path{main.tex},
\path{AGENTS.md}, \path{CODEX.md}, the complete PNP skill and its experiment
mode/protocol, all TASK-001--TASK-007 specifications, actual
\path{main4.tex}/\path{main5.tex}, newest verified MAIN SYNTHESIS-0026, and
newest peers DENSNOGO1, BBTEST1, LOGCOF1, and PHASEBRIDGE1. It inspected
relevant historical MAIN snapshots; PADQUOT1, COVCOLL1, SATQUOT1, WRAPORBIT1,
QNORM1, and earlier padding/trace reports; the HEC, VSC-HEC, G-PPC-HEC,
CRS-PPC, CCSPC, RCCN, and Search-G branches; corpus and predecessor governance;
and the cited experiment/result artifacts. Repository search covered density,
mutation invariance, semantic quotients, padding, alternate representations,
collision entropy, property testing, wrappers, and prior formulations.

\section{DUPLICATION / CONTRADICTION AUDIT}
DENSNOGO1 already proved that one surgically mutable exception has uniform
canonical-layer density \(2^{-k}\). BBTEST1 already proved finite black-box
non-certification. PADQUOT1 already showed that full unpointed padding
invariance collapses on satisfiable independent products. COVCOLL1 already
showed that full program/source covariance collapses to the decision bit. This
report does not re-claim those results. Its new contribution is the unequal-
fiber theorem (5)--(7): even a semantic, representation-invariant bad set has
arbitrarily encoding-sensitive \emph{raw-uniform} mass. There is no
contradiction with canonical-encoding density because that narrower repair is
explicitly retained.

\section{WHAT CHANGED RELATIVE TO ROOT MAIN}
The pre-run LSPV-2 region asked for a density theorem under an explicit
distribution and required encoding normalization. This run adds a missing
condition: the measure itself must be covariant under admitted encodings, or
fiber multiplicities must be bounded. Merely normalizing before decision does
not normalize the sampling distribution. Raw-uniform encoding density is now
retired as a semantics-only bridge; canonical/pushforward density and the
direct whole-execution target remain open.

\section{ADOPTED / REJECTED / RETIRED PRIOR COMPONENTS}
\begin{longtable}{p{0.34\linewidth}p{0.18\linewidth}p{0.38\linewidth}}
\toprule Component & Status & Disposition\\ \midrule
DENSNOGO1 sparse exception & \EXACT & Adopted; raw density is already small.\\
BBTEST1 finite audit collision & \EXACT & Adopted; finite sampling is incomplete.\\
PADQUOT1 full-padding collapse & \EXACT & Adopted; unpointed full invariance erases source identity.\\
Semantic pushforward distribution & \CAND & Smallest noncollapsed density repair.\\
Uniform raw-encoding density & \RET & Arbitrarily changed by fiber multiplicity.\\
Density-to-time transfer & \OPEN & No same-realization theorem.\\
Direct clean-core lower bound & \OPEN & Separation-strength.\\
\bottomrule
\end{longtable}

\section{STALE OR OVERSTATED CLAIMS FOUND}
It is overstated to call a distribution ``uniform'' without naming the sample
space: uniform semantic objects, uniform canonical encodings, and uniform all
valid redundant encodings induce different measures. It is also overstated
to say that parsing/canonicalization alone repairs representation dependence;
the sampler's pushforward must be controlled. No stale theorem in actual
\path{main4.tex} or \path{main5.tex} was promoted; both remain conditional
manuscripts and neither proves an arbitrary-decider density bridge.

\section{Strict status and handoff}
\begin{center}
\begin{tabularx}{0.98\linewidth}{Xp{0.22\linewidth}X}
\toprule Item & Status & Reason\\ \midrule
Redundant-encoding density dilution & \EXACT & Fiber counting\\
Finite encoding/wrapper audit & \COMP & 61,487 bounded assertions\\
Raw-encoding density as universal bridge & \RET & Unequal fibers\\
Canonical or semantic-pushforward density & \OPEN & No all-decider mass theorem\\
Density-to-total-cost transfer & \OPEN & No model-matched implication\\
Direct whole-execution lower bound & \OPEN & Equivalent to separation target\\
P versus NP & \OPEN & Load-bearing arrows remain open\\
\bottomrule
\end{tabularx}
\end{center}

\textbf{Highest-value next target.} Fix one canonical clean-core encoding and
one efficiently samplable semantic distribution. For a precisely defined
provenance-free, wrapper-invariant violation predicate, prove or falsify
inverse-polynomial pushforward mass for every total uniform polynomial SAT
decider. Then still prove a charged \same{} transfer from detection to an
unconditional ordinary-machine lower bound. Otherwise pursue the direct
QNORM1 whole-execution target.

\begin{thebibliography}{99}
\bibitem{GGR98} O. Goldreich, S. Goldwasser, and D. Ron, Property testing and
its connection to learning and approximation, \emph{JACM} 45(4), 653--750,
1998, \url{https://doi.org/10.1145/285055.285060}.
\bibitem{GR23} O. Goldreich and D. Ron, Testing distributions of huge objects,
\emph{TheoretiCS} 2, 2023,
\url{https://doi.org/10.46298/theoretics.23.12}.
\bibitem{FPR26} R. Ferreira Pinto Jr., D. Palit, and S. Raskhodnikova,
Computational Complexity in Property Testing, \emph{SODA 2026},
\url{https://doi.org/10.1137/1.9781611978971.19}; ECCC TR25-185.
\bibitem{BGS75} T. Baker, J. Gill, and R. Solovay, Relativizations of the
\(P=?NP\) question, \emph{SIAM J. Comput.} 4(4), 431--442, 1975,
\url{https://doi.org/10.1137/0204037}.
\bibitem{AW09} S. Aaronson and A. Wigderson, Algebrization: a new barrier in
complexity theory, \emph{ACM TOCT} 1(1), 2009,
\url{https://doi.org/10.1145/1490270.1490272}.
\end{thebibliography}

\bigskip
\noindent\textbf{Final status: P versus NP remains OPEN.}

% END_OF_VERIFIED_REPORT::LSPV-2::TAGDILUTE1
\end{document}
